\documentclass[11pt]{article}
\usepackage[margin=0.9in]{geometry}
\usepackage{amsmath,amssymb,amsthm,booktabs,microtype,xcolor}
\usepackage[hidelinks]{hyperref}
\setlength{\parindent}{0pt}
\setlength{\parskip}{0.55em}
\newtheorem{theorem}{Theorem}
\newtheorem{lemma}[theorem]{Lemma}
\newtheorem{corollary}[theorem]{Corollary}
\newcommand{\B}{\mathcal B}
\newcommand{\C}{\mathbb C}
\newcommand{\Hh}{\mathcal H}
\newcommand{\Ran}{\operatorname{Ran}}
\newcommand{\status}{\textcolor{blue!60!black}{\textbf{candidate partial result---likely valid}}}

\begin{document}

\begin{center}
{\Large\bfseries Rank-one irreducible perturbations from cyclic approximate eigenvectors}\par
\vspace{0.35em}
{\large Partial-result packet for arXiv:2504.17190}\par
\vspace{0.45em}
\status
\end{center}

\section*{1. Source problem and the unresolved bottleneck}

For a separable infinite-dimensional complex Hilbert space $\Hh$, Fang--Jiang--Ma--Shen--Shi--Wang ask:

\begin{quote}
For every $T\in\B(\Hh)$ and $\varepsilon>0$, is there a trace-class $K$ with
$\lVert K\rVert_1<\varepsilon$ such that $T+K$ is irreducible?
\end{quote}

This is \textbf{Problem A} in arXiv:2504.17190v6.  Their Main Theorem reduces it to their type-$\mathrm{II}_1$ factor conjecture: a generator of a type-$\mathrm{II}_1$ factor should admit an arbitrarily small trace-class perturbation that is a direct sum of at most countably many irreducible operators.  The source proves many large subcases but leaves this endpoint open.

The result below gives a different, rank-one sufficient condition.  It does not solve the type-$\mathrm{II}_1$ conjecture, because cyclicity for the von Neumann algebra $W^*(T)$ is much weaker than one-sided cyclicity for $T$ or $T^*$.

We use $u\otimes v$ for the rank-one operator $x\mapsto\langle x,v\rangle u$.  Let $\Omega_\infty(T)$ be the unbounded component of the resolvent set $\rho(T)$.

\section*{2. Main partial theorem}

\begin{theorem}[cyclic approximate-eigenvector criterion]\label{thm:main}
Let $T\in\B(\Hh)$.  Suppose that some
$\lambda\in\partial\Omega_\infty(T)\cap\sigma(T)$ admits unit vectors $\xi_j$ such that
\begin{enumerate}
\item every $\xi_j$ is cyclic for $T^*$, and
\item $\lVert(T-\lambda)\xi_j\rVert\to0$.
\end{enumerate}
Then for every $\varepsilon>0$ there is a rank-one operator $K$ with
$\lVert K\rVert_1<\varepsilon$ such that $T+K$ is irreducible.
\end{theorem}

\begin{corollary}\label{cor:dense}
If the cyclic vectors for $T^*$ are dense in $\Hh$, then $T$ satisfies the conclusion of Theorem~\ref{thm:main}.  The same conclusion holds if the cyclic vectors for $T$ are dense.
\end{corollary}

The second assertion follows by applying the first assertion to $T^*$ and taking adjoints.  Thus this provides a full trace-class-density answer for every operator in either of these co-cyclic classes, and the perturbation is in fact rank one.

\section*{3. Proof}

Fix $\varepsilon>0$.  Choose a unit cyclic vector $\xi$ for $T^*$ with
$\lVert(T-\lambda)\xi\rVert$ as small as needed below.  On $\Omega_\infty(T)$ consider
\[
 g_\xi(z)=\big\langle (T-z)^{-1}\xi,\xi\big\rangle .
\]
This analytic function is not identically zero, since
$g_\xi(z)=-z^{-1}+O(|z|^{-2})$ as $z\to\infty$.  Its zeros are therefore isolated.  We may choose
$\mu\in\Omega_\infty(T)$ arbitrarily close to $\lambda$ such that
\[
 g_\xi(\mu)\ne0,
 \qquad \lVert(T-\mu)\xi\rVert<\varepsilon.
\]
Put $P=\xi\otimes\xi$ and
\[
 K=(\mu I-T)P,
 \qquad S=T+K.
\]
Then $K$ has rank one and
$\lVert K\rVert_1=\lVert(T-\mu)\xi\rVert<\varepsilon$.  Moreover,
\begin{equation}\label{eq:factor}
 S-\mu=(T-\mu)(I-P).
\end{equation}
Consequently $\ker(S-\mu)=\C\xi$ and
$\Ran(S-\mu)=(T-\mu)\xi^\perp$.  The choice $g_\xi(\mu)\ne0$ says exactly that
\[
 (T-\mu)^{-1}\xi\notin\xi^\perp,
\]
or equivalently $\xi\notin\Ran(S-\mu)$.  Thus the eigenvalue $\mu$ is semisimple with one-dimensional eigenspace.

It is also isolated.  Indeed, on $\Omega_\infty(T)$ the possible new spectrum of the rank-one perturbation is the zero set of the analytic Fredholm determinant
\[
 d(z)=\det\bigl(I+(T-z)^{-1}K\bigr).
\]
This determinant tends to $1$ at infinity and hence is not identically zero; its zeros are isolated.  Equation~\eqref{eq:factor} shows $d(\mu)=0$.  Therefore the Riesz projection $E$ of $S$ at $\mu$ has rank one and range $\C\xi$.

The Riesz projection belongs to $W^*(S)$.  Hence the orthogonal range projection
\[
 P=s(EE^*)
\]
also belongs to $W^*(S)$.  On the other hand, $K^*$ has range contained in $\C\xi$.  Induction gives, for every $n\geq0$,
\[
 \operatorname{span}\{\xi,S^*\xi,\ldots,(S^*)^n\xi\}
 =
 \operatorname{span}\{\xi,T^*\xi,\ldots,(T^*)^n\xi\}.
\]
Thus $\xi$ is cyclic for $S^*$.  Finally, for all $m,n\geq0$,
\[
 (S^*)^m P S^n=(S^*)^m\xi\otimes(S^*)^n\xi\in W^*(S).
\]
Because the $S^*$-orbit of $\xi$ is dense, these rank-one operators generate
$\B(\Hh)$ as a von Neumann algebra.  Hence $W^*(S)=\B(\Hh)$, so $S=T+K$ is irreducible.  This proves Theorem~\ref{thm:main}.

For Corollary~\ref{cor:dense}, every point of $\partial\Omega_\infty(T)\cap\sigma(T)$ is an approximate eigenvalue of $T$.  Approximate its unit approximate eigenvectors by cyclic vectors for $T^*$; after normalization the residuals still tend to zero.  The theorem applies.

\section*{4. Relation to the type-$\mathrm{II}_1$ conjecture}

The proof isolates a clean mechanism:
\[
\begin{aligned}
&\text{cyclic approximate eigenvector}\\
&\qquad\Longrightarrow
\text{small rank-one perturbation with a recoverable simple eigenprojection}\\
&\qquad\Longrightarrow W^*(T+K)=\B(\Hh).
\end{aligned}
\]
The first implication is quantitative: the trace norm of the perturbation is exactly
$\lVert(T-\mu)\xi\rVert$.

For a cyclic representation of a type-$\mathrm{II}_1$ factor, a vector cyclic for $W^*(T)$ need only be cyclic under words in $T$ and $T^*$; it need not be cyclic under the powers of $T^*$ alone.  This is the precise point at which the rank-one proof stops short of the source conjecture.

\section*{5. Focused upgrade attempts and obstruction audit}

Six upgrade routes were tested after obtaining Theorem~\ref{thm:main}.

\begin{enumerate}
\item \textbf{Cyclic factor representation $\Rightarrow$ one-sided cyclicity.}
This implication fails already in the finite-factor model: on $L^2(M_n)$, left multiplication by a $*$-generator of $M_n$ generates a cyclic factor representation, but every eigenvalue has the right-module multiplicity and the operator need not be cyclic.

\item \textbf{Finite cyclic multiplicity.}
Replacing one simple eigenvalue by a finite cluster can recover a finite-dimensional Riesz corner, but a trace-norm estimate requires a cyclic tuple consisting of simultaneous small-residual vectors.  Von Neumann-algebra cyclicity does not supply such a tuple.

\item \textbf{Countably many rank-one eigenvalue insertions.}
A countable cyclic set always exists, but nuclearity requires the sum of all residual norms to be finite and arbitrarily small.  There is no mechanism forcing a total cyclic set into a shrinking approximate-eigenspace.

\item \textbf{First make $T^*$ cyclic, then apply the theorem.}
Norm-density results for cyclic operators do not provide perturbations small in trace norm.  The same endpoint ideal obstruction reappears, so this route would merely replace Problem A by another trace-class density problem.

\item \textbf{Recover a rank-one projection without inserting an eigenvalue.}
Adding $i\delta(\xi\otimes\xi)$ kills the old factor commutant when $\xi$ is factor-cyclic, but the projection is not in general recoverable from the single perturbed operator.  The eigenvalue construction solves recovery only at the cost of demanding one-sided cyclicity.

\item \textbf{Relative-normalizer completion.}
The source's averaging lemma succeeds when a diffuse subalgebra of $W^*(\operatorname{Re}T)$ has enough relative normalizers.  Singular masas and strongly solid factors show why such normalizers cannot be assumed in the remaining type-$\mathrm{II}_1$ case.
\end{enumerate}

The remaining gap is therefore structural, not an unverified estimate in the rank-one proof.  A full upgrade would require either a new way to recover finite-rank data from $T+K$ without one-sided cyclicity, or a trace-class theorem producing a cyclic approximate-eigenvector configuration for arbitrary type-$\mathrm{II}_1$ generators.

\section*{6. Status and novelty boundary}

The theorem is a self-contained candidate partial result.  Targeted searches of the cheap run indexes and arXiv for the combinations ``rank-one perturbation'', ``irreducible operator'', and ``cyclic adjoint'' found no exact match, but this is not a comprehensive literature proof of novelty.  Human review should check classical rank-one perturbation and cyclic-operator literature before making a novelty claim.

The source's full Problem A and its equivalent type-$\mathrm{II}_1$ conjecture remain unresolved by this packet.

\end{document}
